\section{Multi-scale realizations of AEC: from cells to organisms}
\label{sec:multiscale}

The AEC definition (Section~\ref{sec:life_aec}) is an interface statement.
This section maps it to familiar biological mechanisms as concrete realizations of information acquisition, predictive modeling, and feedback correction.

\subsection{Cellular scale: homeostasis as error-correction budget management}
\label{sec:cell_scale}

At the cellular scale, the three elements of Definition~\ref{def:agent_aec} have direct correspondences:
\begin{itemize}
  \item \textbf{Information acquisition.}
  Receptor signaling, metabolic sensing, and damage sensing are finite-resolution readouts with thresholds, noise, and discretization.
  They produce conditional information about external conditions and internal state \cite{Alberts2014}.
  \item \textbf{Internal model.}
  Gene-regulatory networks and epigenetic states encode compressed priors mapping environments to phenotypes; operationally, they approximate future readout distributions under finite resources \cite{Alberts2014}.
  \item \textbf{Feedback correction.}
  Proteostasis (folding quality control, chaperones, ubiquitin--proteasome degradation), DNA replication proofreading and repair, membrane-potential maintenance, and stress responses are all mechanisms that reduce effective mismatch density $\sigma(x,\tau)$ by actively exporting dissipation to waste heat and waste material \cite{Alberts2014,Lindahl1993,Hopfield1974,Ninio1975}.
\end{itemize}
The HPA--$\Omega$ framing emphasizes that these processes are not optional ``luxuries''; they are forced by the inevitability of phase friction at finite resolution.

\subsection{Organism scale: behavior and nervous systems as higher-level reverse compilers}
\label{sec:organism_scale}

At the organism scale, reverse compilation (Section~\ref{sec:reverse_compilation}) appears as predictive perception and action:
\begin{itemize}
  \item \textbf{Sensorimotor loops.}
  Sensory streams provide conditional information, internal circuits encode predictive models of future sensory outcomes, and actions implement feedback that shapes future readout \cite{RaoBallard1999,Friston2010}.
  \item \textbf{Learning and memory.}
  Learning increases predictive mutual information rate $\dot I_{\mathrm{pred}}$ by compressing regularities, while memory maintenance and updating incur geometric Landauer costs that depend on circuit architecture (Eq.~\eqref{eq:geom_landauer}) \cite{Friston2010}.
  \item \textbf{Behavioral teleology.}
  Strategy selection can be reinterpreted as maximizing long-run predictive efficiency $\eta_{\mathrm{pred}}$ (Eq.~\eqref{eq:eta_pred}) under metabolic and architectural constraints.
\end{itemize}
This does not replace evolutionary fitness; it provides a physical decomposition of fitness into an information-rate budget and a dissipation budget, yielding a measurable interface between biology and readout thermodynamics.


